% LTeX: language=de

\subsection{Grundlagen}

\begin{Example}{500 x Würfeln}{}
    Ein Würfel wird 500 mal geworfen. Es wird notiert, wie oft eine Augenzahl auftritt.\\[1.0ex]
    $\Omega = \left\{ 1; 2; 3; 4; 5; 6 \right\}$\\[1.5ex]
    \renewcommand{\arraystretch}{1.2}
    \begin{tabularx}{\textwidth}{@{}l|Y|Y|Y|Y|Y|Y@{}}
        Elementarereignis & $\left\{ 1 \right\}$ & $\left\{ 2 \right\}$ & $\left\{ 3 \right\}$ & $\left\{ 4 \right\}$ & $\left\{ 5 \right\}$ & $\left\{ 6 \right\}$ \\
        \hline
        absolute Häufigkeit & 96 & 60 & 90 & 78 & 94 & 82 \\
        \hline
        \multirow{2}{*}{relative Häufigkeit} & $\frac{96}{500}$ & $\frac{60}{500}$ & $\frac{90}{500}$ & $\frac{78}{500}$ & $\frac{94}{500}$ & $\frac{82}{500}$ \\[0.5ex]
        & $= \SI{0.192}{}$ & $= \SI{0.12}{}$ & $= \SI{0.18}{}$ & $= \SI{0.156}{}$ & $= \SI{0.188}{}$ & $= \SI{0.164}{}$ \\
    \end{tabularx}
\end{Example}

\begin{note}
    \texthl{\textbf{Vierfeldertafel}}: stellt den Zusammenhang \texthl{\textbf{zweier} Ereignisse} $A$ und $B$ her
    \begin{tightcenter}
    \renewcommand{\arraystretch}{1.3}
    \begin{tabularx}{0.3\linewidth}[t]{@{} Y|Y|Y @{}}
        & $A$ & $\overline{A}$\\ \hline
        $B$ & $A \cap B$ & $\overline{A} \cap B$ \\ \hline
        $\overline{B}$ & $A \cap \overline{B}$ & $\overline{A} \cap \overline{B}$
    \end{tabularx}
    \end{tightcenter}
\end{note}

\begin{Exercise}{S. 164/2}{}
\begin{enumerate}[leftmargin=*]
    \item [2.] {
        A: bestitzt Autoführerschein;\quad M: besitzt Motorradführerschein\\[1.0ex]
        \begin{minipage}[t]{0.35\linewidth}
            4-Felder Tafel mit absoluten Häufigkeiten:
        \end{minipage}
        \hfill
        \begin{minipage}[t]{0.6\linewidth}
        \renewcommand{\arraystretch}{1.2}
        \begin{tabularx}{0.8\linewidth}[t]{@{} Y|Y|Y|Y @{}}
            $\Omega$ & $A$ & $\overline{A}$ & $\Sigma$ \\ \hline
            $M$ & \cellcolor{mylightgray}$\SI{23}{}$ & $\SI{32}{}$ & \cellcolor{mylightgray} $\SI{55}{}$\\ \hline
            $\overline{M}$ & $\SI{157}{}$ & $\SI{38}{}$ & $\SI{195}{}$ \\ \hline
            $\Sigma$ & \cellcolor{mylightgray}$\SI{180}{}$ & $\SI{70}{}$ & \cellcolor{mylightorange}\tikzmarknode{sum}{$\SI{250}{}$}
        \end{tabularx}

        \begin{tikzpicture}[remember picture, overlay,font=\footnotesize, color=myg]

            \node[align=left, text width=3.2cm] 
                (label1) at ([xshift=2.5cm, yshift=0.0cm] sum.east) {
                    \footnotesize Gesamtzahl
            };

            \draw[-latex]
                ([anchor=west, xshift=0cm]label1) 
                to[out=180, in=0, looseness=1.2]
                ([anchor=east, yshift=0mm, xshift=0.4cm]sum);

        \end{tikzpicture}
        \end{minipage}\\[1.0ex]
        $\left|\overline{A} \cap \overline{M}\right| = 38$\qquad \qquad $\left|A \cup M\right| = \left|\overline{\overline{A} \cap \overline{M}}\right| = 212$
        }
\end{enumerate}
\end{Exercise}

\begin{Exercise}{S. 164/3}{}
\begin{enumerate}[leftmargin=*]
    \item [3.]\begin{enumerate}[label=\alph*),leftmargin=*]
        \item[a)] {
            \begin{minipage}[t]{0.35\linewidth}
            4-Felder Tafel mit relativen Häufigkeiten/Wahrscheinlichkeiten:
            \end{minipage}
            \hfill
            \begin{minipage}[t]{0.6\linewidth}
            \renewcommand{\arraystretch}{1.2}
            \begin{tabularx}{0.8\linewidth}[t]{@{} Y|Y|Y|Y @{}}
                $\Omega$ & $D$ & $\overline{D}$ & $\Sigma$ \\ \hline
                $F$ & \cellcolor{mylightgray}$\frac{4}{50}$ & $\frac{17}{50}$ & \cellcolor{mylightgray} $\frac{21}{50}$\\ \hline
                $\overline{F}$ & $\frac{24}{50}$ & $\frac{5}{50}$ & $\frac{29}{50}$ \\ \hline
                $\Sigma$ & \cellcolor{mylightgray}$\frac{28}{50}$ & $\frac{22}{50}$ & \cellcolor{mylightorange}\tikzmarknode{sum}{$\frac{50}{50}$}
            \end{tabularx}
            \end{minipage}
            }
        \item [b)] {
            $\left|\overline{F} \cap \overline{D}\right| = 50$
            }
        \item [c)] {
            $\left|\overline{F} \cap \overline{D}\right| = \left|F \cup D\right| = 450$ \qquad $P(F \cup D) = \frac{450}{500} = \SI{90}{\%}$
            }
    \end{enumerate}
\end{enumerate}
\end{Exercise}